% Generated by the blueprint tool (blueprint-pdf/1) from a blueprint snapshot; edit the
% Markdown sources, not this file.  Compile with `latexmk -pdf` or by running
% `pdflatex` twice (the second run settles the contents and the page references).
% Images, if any, are read from assets/ next to this file.
\documentclass[11pt]{article}
\usepackage[T1]{fontenc}
\usepackage[utf8]{inputenc}
\usepackage{lmodern}
\usepackage{textcomp}
\usepackage[a4paper,margin=2.5cm,headheight=14pt]{geometry}
\usepackage{amsmath,amssymb,amsthm}
\usepackage{mathtools}
\usepackage{bm}
\usepackage{graphicx}
\usepackage{xcolor}
\usepackage{array,booktabs,tabularx}
\usepackage{fancyvrb}
\usepackage[normalem]{ulem}
\usepackage{microtype}
% Optional, and only with their fonts: a small TeX installation has the
% packages without the fonts, which fails at the first formula.
\IfFileExists{mathrsfs.sty}{\IfFileExists{rsfs10.tfm}{\usepackage{mathrsfs}}{}}{}
\IfFileExists{stmaryrd.sty}{\IfFileExists{stmary10.tfm}{\usepackage{stmaryrd}}{}}{}
\IfFileExists{cancel.sty}{\usepackage{cancel}}{}
\usepackage[hyphens]{url}
\usepackage[pdfusetitle,bookmarksnumbered=false]{hyperref}
\usepackage{bookmark}
\definecolor{bplink}{HTML}{1F4E9C}
\definecolor{bpmuted}{HTML}{5F676E}
\hypersetup{colorlinks=true,linkcolor=bplink,urlcolor=bplink,citecolor=bplink}
\setlength{\parindent}{0pt}
\setlength{\parskip}{0.55em plus 0.2em minus 0.1em}
% Lists as tight as the page's.
\makeatletter
\def\@listi{\leftmargin\leftmargini\parsep 0pt\topsep 0.3em plus 0.1em\itemsep 0.1em plus 0.1em}
\let\@listI\@listi
\def\@listii{\leftmargin\leftmarginii\labelwidth\leftmarginii\advance\labelwidth-\labelsep\parsep 0pt\topsep 0.1em\itemsep 0.1em}
% An ordered list that does not start at 1.
\newcommand{\bpstartat}[1]{\setcounter{\@enumctr}{\numexpr#1-1\relax}}
\makeatother
\setcounter{tocdepth}{3}
\emergencystretch=3em
\allowdisplaybreaks

% A numbered statement, set as a paper sets one: "Theorem 2.1 (Title)."
% #1 is the lead ("Theorem~2.1"), #2 what follows it in roman.
\newenvironment{bpstatement}[2]{\par\addvspace{0.6\baselineskip}\noindent\textbf{#1}#2\textbf{.}\enspace\ignorespaces}
  {\par\addvspace{0.4\baselineskip}}
\newenvironment{bpstatementit}[2]{\begin{bpstatement}{#1}{#2}\itshape\ignorespaces}{\end{bpstatement}}
% A step: a relation with prose of its own, read under the object it starts from.
\newenvironment{bpstep}[1]{\par\addvspace{0.3\baselineskip}\begin{list}{}{\leftmargin=1.6em\rightmargin=0pt\topsep=0pt}\item\relax
  {\small\color{bpmuted}#1}\par\nopagebreak}{\end{list}\par}
% Notes about where an object is written out.
\newenvironment{bpnote}{\par\small\color{bpmuted}\itshape}{\par}
\newcommand{\bpomitted}[1]{\par{\small\color{bpmuted}\itshape #1\par}}
\newcommand{\bpbodyheading}[1]{\par\addvspace{0.4\baselineskip}\noindent\textbf{#1}\par\nopagebreak}
\newcommand{\bprule}{\par\medskip\noindent\hfil\rule{0.3\linewidth}{0.4pt}\hfil\par\medskip}
\newcommand{\bpboundary}[1]{\par\nopagebreak{\small\color{bpmuted}#1\par}\nopagebreak}
% A reference to something this export leaves out, and one that resolves nowhere.
\newcommand{\bpout}[1]{#1\textsuperscript{\hyperlink{bp-outside}{\ensuremath{\ast}}}}
\newcommand{\bpbroken}[1]{\textcolor{red!70!black}{[#1?]}}
\newcommand{\bpslug}[1]{\texttt{#1}}
% A character the export knows no way to set.
\newcommand{\bpglyph}[1]{\fbox{\scriptsize #1}}
% Images, at their own size unless they are wider than the line.
\newsavebox{\bpimagebox}
\newcommand{\bpimage}[1]{\sbox{\bpimagebox}{\includegraphics{#1}}%
  \ifdim\wd\bpimagebox>\linewidth\resizebox{\linewidth}{!}{\usebox{\bpimagebox}}\else\usebox{\bpimagebox}\fi}
\newcommand{\bpmissingimage}[1]{\fbox{\parbox{0.8\linewidth}{\centering\small\color{bpmuted}Image not included: #1}}}
\DefineVerbatimEnvironment{bpverbatim}{Verbatim}{fontsize=\small,xleftmargin=1em}
\renewcommand{\qedsymbol}{\ensuremath{\blacksquare}}
\hypersetup{pdftitle={Compactness through finite intersections — 3 Products},pdfsubject={Section 3 Products; all levels},pdfcreator={blueprint}}
% KaTeX commands the website knows
\providecommand{\R}{\mathbb{R}}\providecommand{\N}{\mathbb{N}}\providecommand{\Z}{\mathbb{Z}}
\providecommand{\Q}{\mathbb{Q}}\providecommand{\C}{\mathbb{C}}
\providecommand{\Reals}{\mathbb{R}}\providecommand{\reals}{\mathbb{R}}\providecommand{\Complex}{\mathbb{C}}
\providecommand{\cnums}{\mathbb{C}}\providecommand{\natnums}{\mathbb{N}}\providecommand{\Bbb}[1]{\mathbb{#1}}
\providecommand{\bold}[1]{\mathbf{#1}}\providecommand{\frak}[1]{\mathfrak{#1}}\providecommand{\mathbbm}[1]{\mathbb{#1}}
\providecommand{\lang}{\langle}\providecommand{\rang}{\rangle}\providecommand{\bra}[1]{\langle #1\rvert}
\providecommand{\ket}[1]{\lvert #1\rangle}\providecommand{\braket}[1]{\langle #1\rangle}
\providecommand{\set}[1]{\{#1\}}\providecommand{\Set}[1]{\left\{#1\right\}}
\providecommand{\infin}{\infty}\providecommand{\isin}{\in}\providecommand{\plusmn}{\pm}\providecommand{\sub}{\subset}
\providecommand{\sube}{\subseteq}\providecommand{\supe}{\supseteq}\providecommand{\exist}{\exists}
\providecommand{\bull}{\bullet}\providecommand{\larr}{\leftarrow}\providecommand{\rarr}{\rightarrow}
\providecommand{\lrarr}{\leftrightarrow}\providecommand{\Larr}{\Leftarrow}\providecommand{\Rarr}{\Rightarrow}
\providecommand{\Lrarr}{\Leftrightarrow}\providecommand{\uarr}{\uparrow}\providecommand{\darr}{\downarrow}
\providecommand{\harr}{\leftrightarrow}\providecommand{\lArr}{\Leftarrow}\providecommand{\rArr}{\Rightarrow}
\providecommand{\hArr}{\Leftrightarrow}\providecommand{\alef}{\aleph}\providecommand{\alefsym}{\aleph}
\providecommand{\real}{\Re}\providecommand{\image}{\Im}\providecommand{\weierp}{\wp}
\providecommand{\Alpha}{\mathrm{A}}\providecommand{\Beta}{\mathrm{B}}\providecommand{\Epsilon}{\mathrm{E}}
\providecommand{\Zeta}{\mathrm{Z}}\providecommand{\Eta}{\mathrm{H}}\providecommand{\Iota}{\mathrm{I}}
\providecommand{\Kappa}{\mathrm{K}}\providecommand{\Mu}{\mathrm{M}}\providecommand{\Nu}{\mathrm{N}}
\providecommand{\Omicron}{\mathrm{O}}\providecommand{\Rho}{\mathrm{P}}\providecommand{\Tau}{\mathrm{T}}
\providecommand{\Chi}{\mathrm{X}}\providecommand{\omicron}{o}\providecommand{\gt}{>}\providecommand{\lt}{<}
\providecommand{\And}{\&}\providecommand{\hdashline}{\hline}\providecommand{\operatornamewithlimits}[1]{\operatorname*{#1}}
\providecommand{\llbracket}{[\![}\providecommand{\rrbracket}{]\!]}\providecommand{\cancel}[1]{#1}
\providecommand{\mathscr}{\mathcal}
% the project's macros ([katex.macros] in blueprint.toml)
\providecommand{\Opens}{}\renewcommand{\Opens}[1]{\mathcal{O}(#1)}
\providecommand{\Top}{}\renewcommand{\Top}{\mathbf{Top}}
\providecommand{\cl}{}\renewcommand{\cl}[1]{\overline{#1}}
\providecommand{\fam}{}\renewcommand{\fam}[1]{\mathcal{#1}}

\begin{document}
\pagestyle{plain}
\begin{center}
{\LARGE\bfseries Compactness through finite intersections\par}
\vspace{0.7em}
{\large Section 3: Products\par}
\vspace{0.5em}
{\small\color{bpmuted}All levels \textperiodcentered{} read along the \texttt{refines} order\par}
{\small\color{bpmuted}Snapshot 4c63101a5eb6\par}
\end{center}
\bigskip

% object: products (blueprint/products/_section.md)
\section*{\texorpdfstring{3\quad Products}{3 Products}}
\phantomsection\label{bp1}%
\addcontentsline{toc}{section}{\texorpdfstring{3\quad Products}{3 Products}}
Finite products of compact spaces are compact.  The tube lemma (\hyperref[bp4]{Lemma 2.2.3}) is stated in Section \bpout{2.2}; it is the whole content of the argument.
% object: thm-product (blueprint/products/thm-product.md)
\phantomsection\label{bp2}%
\begin{bpstatementit}{Theorem~3.1}{ (Products of two compact spaces)}
If \(X\) and \(Y\) are compact, so is \(X \times Y\) with the product topology.

\begin{proof}[{Proof}]
Let \(\fam{W}\) be an open cover of \(X \times Y\).  For each \(x \in X\) the slice \(\{x\} \times Y \cong Y\) is compact, so finitely many members of \(\fam{W}\) cover it; call their union \(W_x\).  By \hyperref[bp4]{Lemma 2.2.3} there is an open \(U_x \ni x\) with \(U_x \times Y \subseteq W_x\).

\begin{enumerate}
\item The \(U_x\) cover \(X\), so by \bpout{Definition 2.2.1} finitely many do: \(U_{x_1}, \dots, U_{x_m}\).
\item Then \(X \times Y = \bigcup_{k=1}^{m} U_{x_k} \times Y \subseteq \bigcup_k W_{x_k}\), a union of finitely many members of \(\fam{W}\).
\end{enumerate}

Compare the closed-set form, \bpout{Theorem 2.2.2}, which gives the same result by a less direct route.
\end{proof}
\end{bpstatementit}

% object: rem-infinite (blueprint/products/rem-infinite.md)
\phantomsection\label{bp3}%
\begin{bpstatement}{Remark~3.2}{ (Arbitrary products)}
For an arbitrary family \((X_\alpha)_{\alpha \in A}\) the statement is Tychonoff's theorem, which needs the axiom of choice; the argument above only gives finite products, by induction on the number of factors: \(\prod_{k=1}^{n+1} X_k \cong \bigl(\prod_{k=1}^{n} X_k\bigr) \times X_{n+1}\).
\end{bpstatement}

% object: lem-tube (blueprint/spaces/compact/lem-tube.md)
\phantomsection\label{bp4}%
\begin{bpnote}Written out here; in the full blueprint it is read under 2.2 Compact spaces, outside this export.\end{bpnote}
\begin{bpstatementit}{Lemma~2.2.3}{ (Tube lemma)}
Let \(Y\) be compact, \(x \in X\), and \(W \subseteq X \times Y\) open with \(\{x\} \times Y \subseteq W\).  Then there is an open \(U \ni x\) with \(U \times Y \subseteq W\).
\end{bpstatementit}
% proof: proof/lem-tube
\begin{proof}[{Proof}]\phantomsection\label{bp5}%
Cover \(\{x\} \times Y\) by basic open sets \(U_y \times V_y \subseteq W\).  The \(V_y\) cover \(Y\), so finitely many do, say \(V_{y_1}, \dots, V_{y_n}\), and \(U = U_{y_1} \cap \dots \cap U_{y_n}\) is the open set wanted.
\end{proof}


\bigskip\phantomsection\hypertarget{bp-outside}{}%
\section*{Outside this export}
\addcontentsline{toc}{section}{Outside this export}
References marked \ensuremath{{}^\ast} are to parts of the blueprint this document leaves out:
\begin{description}
\item[Section 2.2] Compact spaces, in 2 Covers and compact spaces (\bpslug{compact})
\item[Definition 2.2.1] Compact space, in 2 Covers and compact spaces (\bpslug{def-{}\allowbreak{}compact})
\item[Theorem 2.2.2] Compactness by closed sets, in 2 Covers and compact spaces (\bpslug{thm-{}\allowbreak{}fip})
\end{description}
\end{document}
